% Part: proof-theory
% Chapter: natural-deduction
% Section: translation-N2i

\documentclass[../../../include/open-logic-section]{subfiles}

\begin{document}

\olfileid{pt}{nat}{tni}

\olsection{\Log{N2} पासून \Log{G2} मध्ये रूपांतर}

\begin{prop}\ollabel{prop:N2-to-G2}
$\Log{N2c}\ (\Log{N2i}) \Proves
\Gamma \Sequent !A$ असेल, तर
$\Log{G2c}\ (\Log{G2i}) + \Cut \Proves \Gamma'
\Sequent \Delta$. येथे~$\Gamma$ मधील
खुणा काढून मिळणारा !!{formula}s चा
बहुसंच~$\Gamma'$ आहे; तसेच~$!A$
हा~$\lfalse$ नसेल तर
$\Delta = \{!A\}$, आणि असेल तर
$\Delta = \emptyset$.
\end{prop}

\begin{proof}
पुन्हा~$\Gamma \Sequent !A$ च्या
!!{proof}~$\delta$ च्या उंचीवर
विगमन करू. ~$\delta$ हे
$x:!A \Sequent !A$ स्वयंसिद्धक
असेल, तर~$\pi$ हे अनुरूप
$!A \Sequent !A$ स्वयंसिद्धक
आहे; किंवा~$!A \ident \lfalse$
असेल तर
$\lfalse \Sequent \ $ स्वयंसिद्धक आहे.

~$\delta$ चा शेवट
अनुमाननियमात होत असेल, तर
पुढील प्रसंग वेगळे पाहू:
\begin{enumerate}
  \item $R$ हा
  मुख्य सूत्र~$!B \lif !C$
  असलेला~\Intro{\lif} नियम असेल,
  आणि~$x:!B$ हे आधारविधानाच्या
  संदर्भात असेल पण निष्कर्षाच्या
  संदर्भात नसेल, तर~$\delta$
  चा शेवट असा होतो:
  \[
  \AxiomC{}
  \RightLabel{$\delta_1$}
  \Deduce$x:!B, \Gamma \fCenter !C$
  \RightLabel{\Intro{\lif}}
  \UnaryInf$\Gamma \fCenter !B \lif !C$
  \DisplayProof
  \]
  विगमन गृहीतकानुसार
  $!B, \Gamma' \Sequent !C$
  ची !!a{proof}~$\pi_1$
  मिळते; $!C \ident \lfalse$
  असेल तर तिचे रूप
  $!B, \Gamma' \Sequent \ $
  असते. त्यावर~\RightR{\lif}
  लावून~$\pi$ मिळते;
  $!C \ident \lfalse$ असताना
  त्याआधी~\RightR{\Weakening}
  लावावे लागते.
  \item $R$ हा~\Intro{\lif}
  असेल, पण~$x:!B$
  आधारविधानाच्या संदर्भात नसेल,
  तर~$\delta$ चा शेवट असा होतो:
  \[
  \AxiomC{}
  \RightLabel{$\delta_1$}
  \Deduce$\Gamma \fCenter !C$
  \RightLabel{\Intro{\lif}}
  \UnaryInf$\Gamma \fCenter !B \lif !C$
  \DisplayProof
  \]
  विगमन गृहीतकानुसार
  निष्पन्न सूत्र असत्य
  नसेल तेव्हा~$\Gamma' \Sequent !C$
  ची !!a{proof}~$\pi_1$ मिळते;
  असत्याच्या प्रसंगी तिचा
  उत्तरपक्ष रिकामा असतो.
  गरज असल्यास आधी
  उजवीकडील दुर्बलीकरणाने
  असत्य उजवीकडे आणा.
  मग~$!B$ साठी
  \LeftR{\Weakening}
  आणि त्यानंतर
  \RightR{\lif} लावून~$\pi$
  मिळते.
  \item $R$ हा~\Elim{\lif}
  असेल, तर~$\delta$ चा
  शेवट असा होतो:
  \[
  \AxiomC{}
  \RightLabel{$\delta_1$}
  \Deduce$\Gamma_1 \fCenter !B \lif !A$
  \AxiomC{}
  \RightLabel{$\delta_2$}
  \Deduce$\Gamma_2 \fCenter !B$
  \RightLabel{\Elim{\lif}}
  \BinaryInf$\Gamma_1, \Gamma_2 \fCenter !A$
  \DisplayProof
  \]
  येथे~$\Gamma = \Gamma_1 \cup \Gamma_2$.
  विगमन गृहीतकानुसार
  !!{proof}s मिळतात:
  $\pi_1$ ही
  $\Gamma_1' \Sequent !B \lif !A$
  ची आणि~$\pi_2$ ही
  $\Gamma_2' \Sequent !B$ ची.
  दुसऱ्या निष्कर्षसूत्राचा असत्य
  हा प्रसंग असेल, तर दुसरी
  सिद्धता मिळवण्यासाठी
  विगमनातून आलेल्या रिकाम्या
  उत्तरपक्षाच्या सिद्धतेवर
  उजवीकडील दुर्बलीकरण लावा.
  पुढीलप्रमाणे !!{proof}~$\pi$
  मिळते:
  \[
  \AxiomC{}
  \RightLabel{$\pi_1$}
  \Deduce$\Gamma_1' \fCenter !B \lif !A$
  \AxiomC{}
  \RightLabel{$\pi_2$}
  \Deduce$\Gamma_2' \fCenter !B$
  \Axiom$!A \fCenter !A$
  \RightLabel{\LeftR{\lif}}
  \BinaryInf$!B \lif !A, \Gamma_2' \fCenter !A$
  \RightLabel{\Cut}
  \BinaryInf$\Gamma_1', \Gamma_2' \fCenter !A$
  \DisplayProof
  \]
  बहुसंच~$\Gamma_1' \cup \Gamma_2'$
  हा $\Gamma' =
  (\Gamma_1 \cup \Gamma_2)'$
  याच्याशी समान असेलच असे नाही.
  उदाहरणार्थ,~$\Gamma_1$
  आणि~$\Gamma_2$ दोन्हींमध्ये
  $x:!C$ असेल, तर
  $\Gamma_1' \cup \Gamma_2'$
  मध्ये~$!C$ च्या दोन प्रती
  असतात; पण
  $(\Gamma_1 \cup \Gamma_2)'$
  मध्ये एकच प्रत असते.
  योग्य~\LeftR{\Contraction}
  अनुमाने जोडून हे दुरुस्त करता येते.

  $!A \ident \lfalse$
  असेल, तर~$\delta_2$
  वरून विगमनाने मिळणारी
  !!a{proof} एकटी रिकामा
  उत्तरपक्ष देत नाही.
  वरच्या रचनेत
  $!A$ च्या ओळख-स्वयंसिद्धकाऐवजी
  रिकामा उत्तरपक्ष असलेले
  असत्य-स्वयंसिद्धक वापरा.
  डाव्या अभिव्यंजन-नियमापुढे
  कट लावल्यावर
  $\Gamma_1', \Gamma_2' \Sequent \ $
  ची !!{proof}~$\pi$ मिळते.
  त्यामुळे यासाठी
  \LeftR{\Weakening} किंवा
  \RightR{\Weakening}
  लावण्याची गरज नाही.
\item $R$ हा~\FalseCl असेल,
तर~$\delta$ चा शेवट असा:
  \[
  \AxiomC{}
  \RightLabel{$\delta_1$}
  \Deduce$x: \lnot !A, \Gamma \fCenter \lfalse$
  \RightLabel{\FalseCl}
  \UnaryInf$\Gamma \fCenter !A$
  \DisplayProof
  \]
  विगमन गृहीतकानुसार
  $\lnot !A, \Gamma' \Sequent \ $
  ची !!a{proof}~$\pi_1$
  आहे. पुढील !!{proof}~$\pi$ घेऊ:
  \[
  \Axiom$!A \fCenter !A$
  \RightLabel{\RightR{\lnot}}
  \UnaryInf$\fCenter !A, \lnot !A$
  \AxiomC{}
  \RightLabel{$\pi_1$}
  \Deduce$\lnot !A, \Gamma' \fCenter $
  \RightLabel{\Cut}
  \BinaryInf$\Gamma' \fCenter !A$
  \DisplayProof
  \]
\end{enumerate}
फक्त~\FalseCl{} च्या
रूपांतरातून मिळणाऱ्या
!!a{proof}~$\pi$ मध्ये
एखाद्या क्रमवर्तीच्या उजव्या बाजूस
एकापेक्षा अधिक !!{formula}
येऊ शकते. म्हणजे
~\Log{N2i} मधील !!{proof}
रूपांतरित केल्यास
\Log{G2i} मधील !!{proof} मिळते.
\end{proof}

\end{document}
